\documentclass[10pt,conference]{IEEEtran}

% ============================================================
% Packages
% ============================================================
\usepackage[T1]{fontenc}
\usepackage[utf8]{inputenc}
\usepackage{times}
\usepackage{microtype}
\usepackage{graphicx}
\usepackage{booktabs}
\usepackage{array}
\usepackage{tabularx}
\usepackage{enumitem}
\usepackage{xcolor}
\usepackage{hyperref}
\usepackage{minted}
\usepackage{amsmath}
\usepackage{amssymb}

\hypersetup{
    colorlinks=true,
    linkcolor=blue,
    citecolor=blue,
    urlcolor=blue
}

\setlist[itemize]{leftmargin=*, itemsep=2pt, topsep=2pt}
\setlist[enumerate]{leftmargin=*, itemsep=2pt, topsep=2pt}
\setminted{fontsize=\footnotesize,breaklines=true,frame=single,autogobble=true}

% ============================================================

% ============================================================
% Title
% ============================================================
\title{
Assignment 2 Report Template:\\
Domain Adaptation with Small Language Models
}

\author{
\IEEEauthorblockN{
Group Number: \emph{Insert Group Number}
}
\IEEEauthorblockA{
CSCI433/933 Machine Learning Algorithms and Applications\\
Student Names and Student Numbers: \emph{Insert Details}\\
}
}

\begin{document}
\maketitle

% ============================================================
% Page Limit Notice
% ============================================================
\begin{center}
\fbox{
\begin{minipage}{0.94\linewidth}
\textbf{Main Report Page Limit:} The main report must not exceed \textbf{10 pages}, in the style of IEEE conference submissions. 
The 10-page (absolute) limit applies from the abstract through the conclusion.
References, appendices, full evaluation tables, extended prompts, and the GenAI usage log are excluded from the 10-page limit.
Marks are awarded for clarity, technical depth, evidence, and quality of reasoning, not for length.
\end{minipage}
}
\end{center}

% ============================================================
% Abstract
% ============================================================
\begin{abstract}
Briefly summarise the problem, the system developed, the model/RAG approach used, the dataset, the evaluation method, and the main findings. 
This should be approximately 150--200 words.
\end{abstract}

% ============================================================
\section{Introduction and Problem Formulation}
\label{sec:introduction}

\textbf{Relevant marks: Problem formulation and task understanding (10 marks).}

Explain the purpose of the project. 
Describe the task as a domain adaptation problem involving a pretrained Small Language Model (SLM), Shakespearean text, retrieval-augmented generation, and beginner-friendly explanation.

Your introduction should address:

\begin{itemize}
    \item what problem your system is trying to solve;
    \item why Shakespearean text presents a domain adaptation challenge;
    \item who the intended user is;
    \item what your system is expected to do;
    \item what constraints shaped your design.
\end{itemize}

Avoid writing a long generic literature review. 
Focus on the specific problem, system objective, and assignment requirements.

% ============================================================
\section{Dataset and Preprocessing}
\label{sec:dataset}

\textbf{Relevant marks: Dataset preparation and documentation (15 marks).}

Describe the dataset used in the project. 
At minimum, explain how the three required plays were represented and prepared:

\begin{itemize}
    \item \textit{Hamlet};
    \item \textit{Macbeth};
    \item \textit{Romeo and Juliet}.
\end{itemize}

You should describe:

\begin{itemize}
    \item the structure of the dataset;
    \item whether you used utterance-level, scene-level, or other chunking;
    \item metadata retained, such as play, act, scene, speaker, and source identifier;
    \item any cleaning, filtering, or normalisation performed;
    \item any supplementary material used, such as summaries or glossaries;
    \item how primary text was distinguished from supporting explanatory material.
\end{itemize}

\subsection{Dataset Schema}

Provide a concise description of your processed data format. 
For example:

\begin{minted}{json}
{
  "play": "Macbeth",
  "act": 1,
  "scene": 3,
  "speaker": "MACBETH",
  "text": "So foul and fair a day I have not seen.",
  "source_id": "macbeth_1_3_macbeth_001"
}
\end{minted}

\subsection{Chunking Strategy}

Explain your chunking strategy and justify why it is suitable for retrieval. 
Discuss any trade-offs between short chunks, which may be precise but lack context, and longer chunks, which may preserve context but reduce retrieval specificity.

% ============================================================
\section{Methodology}
\label{sec:methodology}

\textbf{Relevant marks: Model adaptation and RAG methodology (20 marks).}

Describe your technical approach. 
This section should explain both your baseline system and your improved RAG-based system.

\subsection{Baseline System}

Describe the baseline system used for comparison. 
This may be a prompt-only model, a simple retrieval-only system, or another minimal approach.

You should explain:

\begin{itemize}
    \item what the baseline does;
    \item what model or tool it uses;
    \item what information it has access to;
    \item why it is a fair baseline for comparison.
\end{itemize}

\subsection{RAG-Based System}

Describe your retrieval-augmented generation pipeline. 
At minimum, explain:

\begin{itemize}
    \item embedding model used;
    \item retrieval/indexing method used;
    \item number of retrieved chunks used for generation;
    \item prompt structure;
    \item language model used for generation;
    \item how retrieved context is incorporated into the generated answer.
\end{itemize}

\subsection{Model Choice and Justification}

Justify your choice of model. 
Discuss how your approach reflects the principles of small, deployable, or resource-constrained language-model systems.

If you used a hosted model, explain why this was appropriate and how the rest of your system design still reflects the assignment's SLM constraints.

% ============================================================
\section{System Design and Implementation}
\label{sec:system}

\textbf{Relevant marks: System implementation and usability (15 marks).}

Describe the implemented system clearly enough that a reader can understand how the components interact.

Include a system diagram if useful.

\subsection{System Pipeline}

Describe the end-to-end flow:

\begin{enumerate}
    \item user query;
    \item query embedding;
    \item retrieval of relevant passages;
    \item prompt construction;
    \item response generation;
    \item evidence display;
    \item logging or evaluation output.
\end{enumerate}

\subsection{Implementation Details}

Summarise the main implementation details:

\begin{itemize}
    \item programming language and libraries;
    \item repository structure;
    \item how to run the system;
    \item dependencies;
    \item any cached embeddings or indices;
    \item limitations of the implementation.
\end{itemize}

Do not paste large blocks of code in the main report. 
Place important snippets only where they support explanation. 
Long code fragments should be placed in the repository, not in the report.

% ============================================================
\section{Evaluation Design}
\label{sec:evaluation-design}

\textbf{Relevant marks: Evaluation, comparison, and error analysis (20 marks).}

Describe your evaluation protocol.

\subsection{Evaluation Questions}

State how many questions were used and distinguish between:

\begin{itemize}
    \item instructor-provided questions;
    \item group-designed questions.
\end{itemize}

Briefly explain what your group-designed questions were intended to test.

\subsection{Scoring Criteria}

Define the scoring scale used for each criterion. 
For example, if you use a 1--5 scale, define what scores 1, 3, and 5 mean.

Your evaluation should include:

\begin{itemize}
    \item correctness;
    \item grounding;
    \item retrieval relevance;
    \item usefulness for a beginner;
    \item style quality, where applicable.
\end{itemize}

\subsection{Evaluation Table Format}

A compact summary table may be included in the main report. 
Full detailed evaluation tables should be placed in the appendix.

\begin{table}[h]
\centering
\caption{Example compact evaluation summary.}
\label{tab:evaluation-summary}
\begin{tabularx}{\linewidth}{lccc}
\toprule
System & Correctness & Grounding & Usefulness \\
\midrule
Baseline & \emph{score} & \emph{score} & \emph{score} \\
RAG System & \emph{score} & \emph{score} & \emph{score} \\
\bottomrule
\end{tabularx}
\end{table}

% ============================================================
\section{Results and Analysis}
\label{sec:results}

\textbf{Relevant marks: Evaluation, comparison, and error analysis (20 marks).}

Present and interpret the results. 
Do not merely list outputs. 
Explain what the results show about the strengths and weaknesses of your system.

\subsection{Baseline versus RAG Comparison}

Compare the baseline system with the improved RAG-based system.

Discuss:

\begin{itemize}
    \item whether RAG improved factual correctness;
    \item whether retrieved evidence improved answer grounding;
    \item cases where RAG helped;
    \item cases where RAG did not help;
    \item any unexpected behaviour.
\end{itemize}

\subsection{Qualitative Examples}

Include a small number of representative examples. 
For each example, show:

\begin{itemize}
    \item the user question;
    \item retrieved evidence;
    \item generated response;
    \item your interpretation of the result.
\end{itemize}

Keep examples concise in the main report. 
Longer examples should be placed in the appendix.

\subsection{Failure Analysis}

Discuss at least three failure cases. 
Possible failure types include:

\begin{itemize}
    \item irrelevant retrieval;
    \item incomplete retrieval;
    \item hallucinated answer;
    \item over-stylised or unclear response;
    \item correct answer but weak evidence;
    \item beginner-unfriendly explanation;
    \item confusion between plays or characters.
\end{itemize}

For each failure case, explain what went wrong and how it might be improved.

% ============================================================
\section{Responsible Use of Generative AI}
\label{sec:genai}

\textbf{Relevant marks: Responsible and transparent GenAI usage (10 marks).}

Summarise how generative AI tools were used in this project.

You should describe:

\begin{itemize}
    \item what tools were used;
    \item what they were used for;
    \item how outputs were verified;
    \item what was accepted, modified, rejected, or improved;
    \item how the group ensured that submitted work reflects its own understanding.
\end{itemize}

The full GenAI usage log may be placed in an appendix or submitted as a separate file.

% ============================================================
\section{Conclusion}
\label{sec:conclusion}

Summarise what was achieved, what was learned, and what limitations remain.

Your conclusion should address:

\begin{itemize}
    \item whether the system met the assignment objectives;
    \item what the evaluation revealed;
    \item what design choices were most important;
    \item what you would improve with more time.
\end{itemize}

% ============================================================
% References
% ============================================================
\begin{thebibliography}{1}

\bibitem{placeholder}
Replace this placeholder with the references used in your report.
Students may alternatively use BibTeX by replacing this environment with:
\texttt{\textbackslash bibliographystyle\{IEEEtran\}} and
\texttt{\textbackslash bibliography\{references\}}.

\end{thebibliography}

% ============================================================
% Appendices
% ============================================================
\appendices

\section{Full Evaluation Table}
\label{app:evaluation}

Include the full evaluation results here.

\begin{table*}[h]
\centering
\caption{Full evaluation table template.}
\label{tab:full-eval}
\begin{tabularx}{\textwidth}{p{2.5cm} p{2.5cm} p{2cm} p{2cm} p{2cm} X}
\toprule
Question & Expected Focus & Correctness & Grounding & Usefulness & Comments \\
\midrule
\emph{Insert question} & \emph{Insert focus} & \emph{Score} & \emph{Score} & \emph{Score} & \emph{Comment} \\
\bottomrule
\end{tabularx}
\end{table*}

\section{Representative Prompts}
\label{app:prompts}

Include important prompts used for retrieval, answer generation, style control, and evaluation.

\begin{minted}{text}
[Insert system prompt or RAG prompt here.]
\end{minted}

\section{Additional Output Examples}
\label{app:outputs}

Include additional examples that support your analysis but do not fit within the 10-page main report.

\section{GenAI Usage Log}
\label{app:genai-log}

The GenAI usage log should include enough information to make the group's use of AI transparent and assessable.

\begin{table*}[h]
\centering
\caption{GenAI usage log template.}
\label{tab:genai-log}
\begin{tabularx}{\textwidth}{p{2.2cm} p{3.0cm} p{3.0cm} X}
\toprule
Tool & Prompt or Task & Useful Output & Verification / Modification \\
\midrule
\emph{Tool name} & \emph{Prompt summary} & \emph{Response summary} & \emph{How it was checked or modified} \\
\bottomrule
\end{tabularx}
\end{table*}

\end{document}
